\section{Цифровая оценка инвестиций в инновационное производство}

Цифровая оценка инвестиций в инновационное производство направлена на
определение экономической эффективности разработки программного
продукта и обоснование целесообразности вложения средств в его
создание. Для этого используются показатели, основанные на анализе
денежных потоков проекта и учёте фактора времени.

В рамках данного раздела для оценки инвестиционной привлекательности
проекта используются следующие показатели:
\begin{itemize}
\item чистая приведённая стоимость ($NPV$);
\item внутренняя норма доходности ($IRR$);
\item рентабельность инвестиций ($ROI$);
\item экономическая добавленная стоимость ($EVA$);
\item срок окупаемости проекта.
\end{itemize}

В качестве исходного объёма инвестиций принимаются затраты на
изготовление продукта:

\begin{equation}
K = 1443246.44 \text{ руб.}
\end{equation}

Согласно результатам предыдущего раздела, годовая прибыль проекта
составляет:

\begin{equation}
C_{\text{год}} = 2178000 \text{ руб.}
\end{equation}

Для расчётов примем, что прибыль поступает равномерно в течение двух
лет, а ставка дисконтирования составляет:

\begin{equation}
r = 0.25.
\end{equation}

\subsection*{Чистая приведённая стоимость}

Чистая приведённая стоимость определяется по формуле:

\begin{equation}
NPV = -K + \sum_{t=1}^{n} \frac{CF_t}{(1+r)^t},
\end{equation}

где $CF_t$ --- денежный поток в период $t$; $n$ --- число расчётных
периодов.

Для рассматриваемого проекта принимаем:
\[
CF_1 = CF_2 = 2178000 \text{ руб.}, \qquad n = 2.
\]

Тогда:

\begin{equation}
NPV = -1443246.44 + \frac{2178000}{1.25} +
\frac{2178000}{1.25^2}.
\end{equation}

После вычисления получаем:

\begin{equation}
NPV = 1693073.56 \text{ руб.}
\end{equation}

Положительное значение $NPV$ показывает, что проект обладает
инвестиционной привлекательностью и обеспечивает прирост стоимости
вложенного капитала.

\subsection*{Внутренняя норма доходности}

Внутренняя норма доходности определяется как такое значение ставки
дисконтирования, при котором чистая приведённая стоимость равна нулю:

\begin{equation}
0 = -K + \sum_{t=1}^{n} \frac{CF_t}{(1+IRR)^t}.
\end{equation}

Для рассматриваемого проекта численным методом получаем:

\begin{equation}
IRR \approx 119.62\%.
\end{equation}

Значение $IRR$ существенно превышает принятую ставку дисконтирования,
что свидетельствует о высокой эффективности инвестиций в разработку
продукта.

\subsection*{Рентабельность инвестиций}

Рентабельность инвестиций определяется как отношение прибыли к объёму
инвестиций:

\begin{equation}
ROI = \frac{C_{2} - K}{K}\cdot 100\%,
\end{equation}

где $C_2$ --- прибыль за два года реализации проекта.

Так как:
\[
C_2 = 4356000 \text{ руб.},
\]

получаем:

\begin{equation}
ROI = \frac{4356000 - 1443246.44}{1443246.44}\cdot 100\%
\approx 201.82\%.
\end{equation}

Следовательно, за два года проект приносит прибыль, более чем в два
раза превышающую первоначальные инвестиции.

\subsection*{Экономическая добавленная стоимость}

Экономическая добавленная стоимость рассчитывается по формуле:

\begin{equation}
EVA = P - WACC \cdot K,
\end{equation}

где $P$ --- чистая прибыль за период; $WACC$ ---
средневзвешенная стоимость капитала.

Примем:
\[
WACC = 0.2.
\]

Тогда при годовой прибыли $P = 2178000$ руб. получаем:

\begin{equation}
EVA = 2178000 - 0.2 \cdot 1443246.44 =
1889350.71 \text{ руб.}
\end{equation}

Положительное значение $EVA$ показывает, что проект создаёт
добавленную стоимость и экономически оправдан с позиции инвестора.

\subsection*{Срок окупаемости}

Срок окупаемости определяется как отношение объёма инвестиций к
среднегодовой прибыли:

\begin{equation}
T_{\text{ок}} = \frac{K}{C_{\text{год}}}.
\end{equation}

Тогда:

\begin{equation}
T_{\text{ок}} = \frac{1443246.44}{2178000} \approx 0.66 \text{ года}.
\end{equation}

В пересчёте на месяцы срок окупаемости составляет:

\begin{equation}
T_{\text{ок}} \approx 7.95 \text{ месяцев}.
\end{equation}

Это означает, что вложения в разработку продукта возвращаются менее
чем за один год эксплуатации.

\subsection*{Учёт затрат по видам деятельности}

Основные методы учёта затрат включают:
\begin{itemize}
\item учёт затрат по видам деятельности;
\item полную стоимость владения.
\end{itemize}

Учёт затрат по видам деятельности представляет собой метод
управленческого учёта, позволяющий оценить затраты на различные виды
деятельности в рамках инновационного процесса. Данный подход даёт
возможность выделить прямые и косвенные расходы, возникающие на этапах
разработки, внедрения и последующей эксплуатации продукта.

Для рассматриваемого проекта в качестве прямых затрат выступают
затраты на заработную плату исполнителей, оборудование,
амортизационные отчисления, расходные материалы и организацию рабочих
мест. Косвенные затраты представлены накладными расходами, связанными
с управлением проектом, сопровождением и обеспечением взаимодействия
между участниками процесса.

Для укрупнённой оценки затрат по видам деятельности используем
показатель полной стоимости владения (\textit{Total Cost of Ownership},
$TCO$), который определяется как сумма затрат на изготовление продукта
и затрат на внедрение всех планируемых экземпляров:

\begin{equation}
TCO = K + N_p^o \cdot K_{\text{вн}},
\end{equation}

где $K$ --- стоимость изготовления продукта;
$N_p^o$ --- планируемое количество реализуемых экземпляров;
$K_{\text{вн}}$ --- стоимость внедрения одного экземпляра.

Согласно предыдущим расчётам:
\[
K = 1443246.44 \text{ руб.},
\qquad
N_p^o = 30,
\qquad
K_{\text{вн}} = 84654.43 \text{ руб.}
\]

Тогда:

\begin{equation}
TCO = 1443246.44 + 30 \cdot 84654.43 = 3982879.34 \text{ руб.}
\end{equation}

Полученное значение показывает полную совокупную стоимость владения
продуктом на горизонте двух лет, включая затраты на его изготовление и
последующее внедрение в планируемом объёме.

\subsection*{Оценка использования инновации на уровне предприятия}

Оценка использования инновации на уровне предприятия позволяет
определить, насколько внедрение разрабатываемого решения способствует
повышению эффективности использования ресурсов организации.

Одним из показателей такой оценки является показатель
производительности инновации $IP$, который определяется как отношение
экономической добавленной стоимости к коммерческим, общим и
административным расходам:

\begin{equation}
IP = \frac{EVA}{SG\&A},
\end{equation}

где $EVA$ --- экономическая добавленная стоимость;
$SG\&A$ --- коммерческие, общие и административные расходы.

В рамках укрупнённой оценки в качестве $SG\&A$ примем суммарные
накладные расходы по проекту, включающие накладные расходы на этапе
изготовления и на этапе внедрения:

\[
SG\&A = C_{\text{накл}} + C_{\text{вн.накл}} =
251428.57 + 830353.50 = 1081782.07 \text{ руб.}
\]

Экономическая добавленная стоимость ранее была рассчитана как:

\[
EVA = 1889350.71 \text{ руб.}
\]

Тогда:

\begin{equation}
IP = \frac{1889350.71}{1081782.07} \approx 1.75
\end{equation}

Поскольку полученное значение превышает единицу, можно сделать вывод,
что внедрение разработанного решения положительно влияет на
эффективность использования ресурсов предприятия. Иными словами,
создаваемая инновация обеспечивает экономический эффект, превышающий
уровень сопутствующих административных и накладных расходов.

\subsection*{Итоговая оценка инвестиционной привлекательности проекта}

Для удобства сведём полученные показатели в таблицу.

\begin{table}[h]
\centering
\caption{Показатели эффективности инвестиций проекта}
\label{tab:investment_metrics}
\begin{tabular}{|p{8cm}|c|}
\hline
Показатель & Значение \\ \hline
Чистая приведённая стоимость $NPV$ & 1693073.56 руб. \\ \hline
Внутренняя норма доходности $IRR$ & 119.62\% \\ \hline
Рентабельность инвестиций $ROI$ & 201.82\% \\ \hline
Экономическая добавленная стоимость $EVA$ & 1889350.71 руб. \\ \hline
Полная стоимость владения $TCO$ & 3982879.34 руб. \\ \hline
Показатель использования инновации $IP$ & 1.75 \\ \hline
Срок окупаемости & 7.95 месяцев \\ \hline
\end{tabular}
\end{table}

Таким образом, проведённая цифровая оценка инвестиций показывает, что
проект является экономически эффективным. Положительное значение
$NPV$, высокий уровень $IRR$, положительная величина $EVA$, значение
показателя $IP$, превышающее единицу, и короткий срок окупаемости
подтверждают целесообразность инвестирования в разработку
рассматриваемого программного продукта.
